\section{Causal tests separate a nonempty zero divisor}
\label{sec:causal-trace-separation}

\paragraph{MA-JP-004 (resolvent-power separation, including two polar moments).}
Let
\[
 D=\sum_{\lambda\in\mathcal Z}d_\lambda[\lambda],
 \qquad d_\lambda\in\mathbb Z_{>0},\qquad
 u\le\operatorname{Re}\lambda\le v,
 \tag{MA-JP.40}
\]
be a locally finite divisor in a fixed vertical strip.
Assume its count includes multiplicities and satisfies
\[
 \sum_{\substack{\lambda\in\mathcal Z\\
                   T\le\operatorname{Im}\lambda<T+1}}
     d_\lambda\le C\log(2+|T|)\qquad(T\in\mathbb R).
 \tag{MA-JP.41}
\]
For $g\in\mathcal D_+=C_c^\infty((0,\infty))$, put
\[
 E_a(g)=\int_0^\infty g(t)e^{at}\,dt,\qquad
 \tau_D(g)=\sum_{\lambda\in\mathcal Z}d_\lambda E_\lambda(g).
 \tag{MA-JP.42}
\]
This sum is absolutely convergent and defines a continuous
linear functional on $\mathcal D_+$.
It is identically zero if and only if $D$ is empty.

More precisely, fix two distinct prescribed complex
exponents $a_1,a_2$.  Then
\[
 \begin{split}
 \tau_D|_{\ker E_{a_1}\cap\ker E_{a_2}}=0
 &\quad\Longleftrightarrow\quad
 \tau_D\in\operatorname{span}\{E_{a_1},E_{a_2}\}\\
 &\quad\Longleftrightarrow\quad
 \mathcal Z\subset\{a_1,a_2\}.
 \end{split}
 \tag{MA-JP.43}
\]
If there is a divisor point outside those two exponents,
therefore, there exists $g\in\mathcal D_+$ with
\[
       E_{a_1}(g)=E_{a_2}(g)=0,
                         \qquad \tau_D(g)=1.
 \tag{MA-JP.44}
\]
The scalar normalization may be complex.
The support is strictly positive; no nonnegativity of
the values of $g$ is asserted.

\begin{proof}
\emph{The functional and its convergence.}
On a fixed compact interval inside $(0,\infty)$,
integration by parts $N$ times gives
\[
 |E_{\sigma+i\gamma}(g)|
  \le C_{N,u,v,\operatorname{supp}g}
        (1+|\gamma|)^{-N}
                \max_{0\le j\le N}\|g^{(j)}\|_\infty
                  \quad(u\le\sigma\le v).
 \tag{MA-JP.45}
\]
Equation (MA-JP.41) makes the sum absolutely convergent
already with $N=2$.  The same estimate proves continuity
on each fixed-support test space and hence on its
usual inductive limit $\mathcal D_+$.

\emph{A compact-test approximation to a resolvent power.}
Choose an integer $k\ge4$ and a real
$\kappa>\max(v,\operatorname{Re}a_1,
\operatorname{Re}a_2)+1$.
Let smooth real profiles $\psi,\chi$ on $[0,\infty)$
have values in $[0,1]$, with
\[
 \begin{aligned}
 &\psi=0\text{ on }[0,1],\quad \psi=1\text{ on }[2,\infty),\\
 &\chi=1\text{ on }[0,1],\quad \chi=0\text{ on }[2,\infty).
 \end{aligned}
\]
For $0<\epsilon<1/4$ and $R\ge1$ set
\[
 g_{\epsilon,R}(t)=t^{k-1}e^{-\kappa t}
                    \psi(t/\epsilon)\chi(t/R).
 \tag{MA-JP.46}
\]
These are actual members of $\mathcal D_+$.
If $\sigma$ lies in any fixed interval containing
$[u,v]$ and $\operatorname{Re}a_1,
\operatorname{Re}a_2$ below $\kappa$, both the
$L^1$ norm and the second-derivative $L^1$ norm of
$e^{\sigma t}g_{\epsilon,R}(t)$ are bounded independently
of $\epsilon,R$.
Indeed the differentiated lower-transition terms are
supported on $[\epsilon,2\epsilon]$; a total of $j\le2$
derivatives there gives an integral $O(\epsilon^{k-j})$.
The terms without a differentiated lower profile have
the integrable powers $t^{k-1-j}$ near zero.
On $[R,2R]$ all differentiated upper profiles are
controlled by a polynomial in $R$ times a fixed
decreasing exponential.  The remaining integrals
are bounded by the same exponential moments.
Thus twice integrating by parts, with zero endpoint
terms for each compact test, proves the uniform bound
\[
 |E_{\sigma+i\gamma}(g_{\epsilon,R})|
                  \le C_\kappa(1+\gamma^2)^{-1}.
 \tag{MA-JP.47}
\]
Pointwise dominated convergence in the integral gives
\[
 E_\lambda(g_{\epsilon,R})\longrightarrow
       \int_0^\infty t^{k-1}e^{-(\kappa-\lambda)t}\,dt
       =\frac{\Gamma(k)}{(\kappa-\lambda)^k}.
 \tag{MA-JP.48}
\]
Equations (MA-JP.41) and (MA-JP.47) justify this
limit in the entire divisor sum, with both endpoint
cutoffs removed.  This justification is needed:
the uncut resolvent kernel itself is not a compact
smooth test supported away from zero.

\emph{Meromorphic pole separation.}
Suppose first that
$\tau_D=A_1E_{a_1}+A_2E_{a_2}$ on $\mathcal D_+$.
Apply that equality to (MA-JP.46), take the justified
limit, and cancel $\Gamma(k)$.  For every sufficiently
large real $\kappa$ this gives
\[
 F_D(\kappa):=
    \sum_{\lambda\in\mathcal Z}
                     \frac{d_\lambda}{(\kappa-\lambda)^k}
      =\frac{A_1}{(\kappa-a_1)^k}
             +\frac{A_2}{(\kappa-a_2)^k}.
 \tag{MA-JP.49}
\]
The series on the left converges normally on every
compact subset avoiding $\mathcal Z$.
For its tail, $\kappa$ in a fixed compact set gives
$|\kappa-\lambda|\ge |\operatorname{Im}\lambda|/2$
once the latter is sufficiently large; (MA-JP.41)
then proves normal convergence.  The remaining
finitely many terms are holomorphic off their poles.
Thus $F_D$ is meromorphic on $\mathbb C$, with
principal part
$d_{\lambda_0}(\kappa-\lambda_0)^{-k}$ at every
$\lambda_0\in\mathcal Z$.
There are no other singular terms there and this
coefficient is nonzero.

The identity theorem first applies on the right
half-plane and then on the complement of the locally
finite discrete union
$\mathcal Z\cup\{a_1,a_2\}$, which is connected.
Consequently (MA-JP.49) is a meromorphic identity.
Any divisor point outside $\{a_1,a_2\}$ would give
a noncancelling pole on the left and a holomorphic
right side there, a contradiction.
If the divisor is supported in those two points,
its defining formula is plainly in their span.
The same argument with a zero right side proves
$\tau_D=0$ on $\mathcal D_+$ if and only if the
divisor is empty.  In fact the coefficients $A_i$
in an identity are exactly the divisor lengths
at $a_i$, with zero when that point is absent.

\emph{The exact linear-algebra step imposing both moments.}
The map
$E=(E_{a_1},E_{a_2}):\mathcal D_+\to\mathbb C^2$
is onto.  To check this concretely, choose a smooth
bump $b$ supported in a sufficiently small interval
about a positive point so that both $E_{a_i}(b)$
are nonzero.  This is possible by approximating
a positive point mass and using continuity of
the two exponential functions.
Choose a positive shift $\Delta$ for which
$e^{(a_1-a_2)\Delta}\ne1$.
The bump and its right translate $b(t-\Delta)$
both belong to $\mathcal D_+$, and their moment
vectors have determinant
\[
 E_{a_1}(b)E_{a_2}(b)
          \bigl(e^{a_2\Delta}-e^{a_1\Delta}\bigr)\ne0.
 \tag{MA-JP.50}
\]
If $\tau_D$ vanishes on $\ker E$, it descends to a
linear functional on the actual two-dimensional
quotient identified by $E$.  It is therefore a
linear combination of $E_{a_1},E_{a_2}$.
The converse is immediate.  This proves the first
equivalence of (MA-JP.43); the pole argument proves
the second.
When the support is not contained in those two
points, a vector in $\ker E$ consequently has
nonzero $\tau_D$ value.  Dividing it by that value
gives (MA-JP.44).
\end{proof}

\paragraph{Application to the actual trace and the fast counterpath.}
Take the actual retained divisor from MA-WP,
$\lambda=\rho-1/2$ with positive length $r_\rho$,
and discard terms with $r_\rho=0$.
It satisfies (MA-JP.40)--(MA-JP.41), since
$0<r_\rho\le m_\rho$ and the fixed-function local
zero count is the one used in MT-M18 and MA-WP.
It is contained in a fixed vertical strip and
has no point at the polar exponents $\pm1/2$.
For these compact tests MA-WP and its actual
ambient lift already prove
\[
 \tau_D(g)=\operatorname{Tr}_{H/J}C_g^{H/J}
                    =\operatorname{Tr}_H(C_gP_Z).
 \tag{MA-JP.51}
\]
Thus the functional separated above is the trace
of an actual trace-class operator, not a scalar
substitute for an unconstructed quotient action.
If the retained divisor is nonempty, (MA-JP.44)
with $a_1=1/2,a_2=-1/2$ supplies the stronger
version of MA-JP-002:
\[
 \begin{gathered}
 g\in C_c^\infty((0,\infty)),\qquad g(0)=0,\\
 G(1/2)=G(-1/2)=0,\qquad
 \operatorname{Tr}_{H/J}C_g^{H/J}=1.
 \end{gathered}
 \tag{MA-JP.52}
\]
The actual counterpath construction MA-JP-003
therefore applies with a causal test, without
any other change to its trace-ideal argument.
For the rational trivial sector, Hardy's established
existence of a critical-line zero (1914, printed
p.~1012; full reference and source scope in MA-JP-003) supplies a
retained point for every $\alpha>0$ and $s\ge0$.
This is an explicit unconditional zero-location
\emph{existence} input, already identified in
MA-JP-003; it is neither RH nor an assumption of
an exterior zero.  For a different function the
nonemptiness hypothesis is retained unless its
own suitable zero has separately been supplied.

\paragraph{Application to the exact slow-band correction.}
Apply the same lemma to the upper divisor
\[
 D_{\eta,+}^{\rm out}
  =\sum_{\operatorname{Re}\rho_\eta\ge1/2+\alpha}
          m_\eta(\rho_\eta)[\rho_\eta-1/2].
 \tag{MA-JP.53}
\]
It satisfies the same count and strip assumptions.
The actual correction in MA-SB-003, or MA-SB-004
for the two character families, vanishes for every
causal compact test if and only if this divisor
is empty.  The equivalence remains true when
both polar moments are required to vanish,
since nontrivial completed zeros do not give
the exponents $\pm1/2$.
If this divisor is nonempty, there is a normalized
causal test for which the actual correction is one.
This last statement is conditional on nonemptiness;
the lemma proves no exterior-zero existence or
zero-location theorem.

The separation proof uses finite-length multiplicities
through actual traces and a normally convergent
meromorphic resolvent-power function.  It neither
requires simple zeros nor assigns a sign to values
of the complex test.  In the trace-framework placement
of BR-TR-01/02, the retained finite-length blocks
contribute their lengths to the trace; the analytic
separation of those contributions is the additional
argument proved here, not an arithmetic positivity
consequence of that categorical description.
